\begingroup
\newcommand{\TransBracket}[1]{\boldsymbol{[}#1\boldsymbol{]}}
\ArithmeticHead{300}{L}{ANALYSIS. NACHLASS.}
\SourceLeaf{307}
\NumberedItem{12}
\noindent[S. 47]\par
\par\indent Zum Beweise der ersten Sätze hingegen wird man sich auf folgende Art verhalten.
\[\begin{aligned}
&(1-xy)(1+xxy)\left(1+\frac{x^8}{y}\right)\left(1-\frac{x^9}{y}\right)\\
&\qquad(1-x^{11}y)(1+x^{12}y)\ \text{etc.}
\end{aligned}\]
\noindent wird entwickelt in
\[\begin{aligned}
&P\left\{\begin{aligned}
&1-x^3yy+x^{26}y^4-x^{69}y^6+{.}{.}{.}\\
&\phantom{1}{-}\frac{x^{17}}{yy}+\frac{x^{54}}{y^4}-\frac{x^{111}}{y^6}+{.}{.}{.}
\end{aligned}\right\}\\[1ex]
&{}+Q\left\{\begin{aligned}
&\phantom{{}-{}}y-x^{13}y^3+x^{46}y^5-{.}{.}{.}\\
&-\frac{x^7}{y}+\frac{x^{34}}{y^3}-\frac{x^{81}}{y^5}+{.}{.}{.}
\end{aligned}\right\}
\end{aligned}\]
\noindent wo $P$ und $Q$ bloss Functionen von $x$ sind. Um diese zu bestimmen, setzen wir erst $y=x^{\frac72}$, wodurch
\[\begin{aligned}
&(1-x^9)(1-x^{11})(1-x^{29})(1-x^{31}){.}{.}{.}\\
&\qquad=P(1-2x^{10}+2x^{40}-{.}{.}{.})
=P\frac{\TransBracket{x^{10}}^2}{\TransBracket{x^{20}}}.
\end{aligned}\]
\noindent Zweitens setzen wir
\[
y=\frac1x,
\]
\noindent so wird
\[\begin{aligned}
&P\{1-x-x^{19}+x^{22}+x^{58}{.}{.}{.}\}\\
&\qquad+\frac Qx\{1-x^9-x^{11}+x^{38}+x^{42}-x^{87}{.}{.}{.}\}=0,\ \text{[*)]}
\end{aligned}\]
\noindent also
\[\begin{aligned}
Q&=-xP\frac{(1-x)(1-x^{19})(1-x^{21}){.}{.}{.}}{(1-x^9)(1-x^{11})(1-x^{29}){.}{.}{.}}\\
&=-x\left[\frac{\TransBracket{x^{20}}}{\TransBracket{x^{10}}^2}\right](1-x)(1-x^{19})(1-x^{21}){.}{.}{.}
\end{aligned}\]
\noindent Setzen wir also
\[\begin{aligned}
(1-x)(1+xx)(1+x^8)(1-x^9){.}{.}{.}&=t\\
(1-x^3)(1+x^4)(1+x^6)(1-x^7){.}{.}{.}&=u,
\end{aligned}\]
\noindent so hat man (aus $y=1$ und $y=xx$)
\Needspace{3\baselineskip}
\par\smallskip\noindent\rule{.25\linewidth}{.4pt}\par
{\small\par\indent [*) In der Handschrift steht $x^{67}$ statt $x^{87}$.]\par}
\clearpage
\ArithmeticHead{301}{R}{{\scriptsize ZUR THEORIE DER TRANSSCENDENTEN FUNCTIONEN. I.}}
\SourceLeaf{308}
\[\begin{aligned}
t&=P(1-x^3-x^{17}+x^{26}+x^{54}-{.}{.}{.})\\
&\qquad+Q(1-x^7-x^{13}+{.}{.}{.})\\[.8ex]
u&=P(1-x^7-x^{13}+x^{34}+x^{46}-{.}{.}{.})\\
&\qquad+Qxx(1-x^3-x^{17}+x^{26}+x^{54}{.}{.}{.}).
\end{aligned}\]
\ClosingRule
\par\indent Man findet aber auch durch ein ganz ähnliches Verfahren
\[\begin{aligned}
&(1-x^3y)(1+x^4y)\left(1+\frac{x^6}{y}\right)\left(1-\frac{x^7}{y}\right){.}{.}{.}\\
&=P\left(1-x^7yy-\frac{x^{13}}{yy}\quad\begin{aligned}
&+x^{34}y^4{.}{.}{.}\\
&+\frac{x^{46}}{y^4}{.}{.}{.}
\end{aligned}\right)\\[1ex]
&\qquad+Qxx\left(\begin{aligned}
&\phantom{{}-{}}y-x^{17}y^3+{.}{.}{.}\\
&-\frac{x^3}{y}+\frac{x^{26}}{y^3}-{.}{.}{.}
\end{aligned}\right).
\end{aligned}\]
\ClosingRule
\par\indent Hieraus folgt
\[\begin{aligned}
tx+u&=(P+Qx)\{1+x-x^4-x^7-x^{13}-x^{18}+x^{27}+x^{34}{.}{.}{.}\}\\
-tx+u&=(P-Qx)\{1-x+x^4-x^7-x^{13}+x^{18}-x^{27}+x^{34}{.}{.}{.}\}
\end{aligned}\]
\noindent oder
\[\begin{aligned}
tx+u&=(P+Qx)\TransBracket{-x^5}(1+x)(1-x^4)(1-x^6)(1+x^9){.}{.}{.}\\
u-tx&=(P-Qx)\TransBracket{-x^5}(1-x)(1+x^4)(1+x^6)(1-x^9){.}{.}{.}
\end{aligned}\]
\noindent Nun ist ferner
\[\begin{aligned}
&(P+Qx)(1-x-x^{19}+x^{22}+{.}{.}{.})\\
&\qquad=-\frac Qx(1-xx+x^3-x^9-x^{11}+x^{21}+x^{24}{.}{.}{.})\\
&\qquad=-\frac Qx\TransBracket{-x^5}(1-xx)(1+x^3)(1+x^7)(1-x^8){.}{.}{.}
\end{aligned}\]
\[\begin{aligned}
&(P-Qx)(1-x-x^{19}+x^{22}+{.}{.}{.})\\
&\qquad=-\frac Qx(1+xx-x^3-x^9-x^{11}-x^{21}+x^{24}\quad{.}{.})\\
&\qquad=-\frac Qx\TransBracket{-x^5}(1+xx)(1-x^3)(1-x^7)(1+x^8){.}{.}{.},
\end{aligned}\]
\noindent also
\[\begin{aligned}
P+Qx&=\frac{\TransBracket{-x^5}}{\TransBracket{x^{10}}^2}(1-xx)(1+x^3)(1+x^7)(1-x^8){.}{.}{.}\\
P-Qx&=\frac{\TransBracket{-x^5}}{\TransBracket{x^{10}}^2}(1+xx)(1-x^3)(1-x^7)(1+x^8){.}{.}{.}
\end{aligned}\]
\[\begin{aligned}
tx+u&=(1+x^5)^2(1+x^{15})^2{.}{.}{.}\\
&\qquad.(1+x)(1-xx)(1+x^3)(1-x^4)\\
&\qquad\phantom{.}(1-x^6)(1+x^7)(1-x^8)(1+x^9){.}{.}{.}
\end{aligned}\]
\clearpage
\ArithmeticHead{302}{L}{ANALYSIS. NACHLASS.}
\SourceLeaf{309}
\[\begin{aligned}
u-tx&=(1+x^5)^2(1+x^{15})^2{.}{.}{.}\\
&\qquad.(1-x)(1+xx)(1-x^3)(1+x^4)\\
&\qquad\phantom{.}(1+x^6)(1-x^7)(1+x^8)(1-x^9){.}{.}{.}
\end{aligned}\]
\noindent[S. 48]\par
\par\indent Hieraus folgt
\[\begin{aligned}
tx+u&=\frac{\TransBracket{xx}^3}{\TransBracket{x}\TransBracket{x^4}}\cdot\frac{\TransBracket{x^{10}}}{\TransBracket{x^5}\TransBracket{x^{20}}}\\
&=\frac{\TransBracket{x}\TransBracket{x^4}}{\TransBracket{xx}^2}\frac{\TransBracket{x^{10}}}{\TransBracket{x^5}\TransBracket{x^{20}}}(1+2x+2x^4+2x^9{.}{.}{.})\\[1ex]
u-tx&=\frac{\TransBracket{x}\TransBracket{x^4}}{\TransBracket{xx}^2}\cdot\frac{\TransBracket{x^{10}}^6}{\TransBracket{x^5}^3\TransBracket{x^{20}}^3}\\
&=\frac{\TransBracket{x}\TransBracket{x^4}}{\TransBracket{xx}^2}\frac{\TransBracket{x^{10}}}{\TransBracket{x^5}\TransBracket{x^{20}}}(1+2x^5+2x^{20}+2x^{45}{.}{.}{.}).
\end{aligned}\]
\noindent Woraus das Theorem [die beiden ersten Gleichungen von art. 10., S. 297, 298] von selbst folgt.
\clearpage
\endgroup
